\chapter{Style}
Now we will be covering some styling options, to make our setup not look like shit. Remember that this may differ depending on your components, so before changing your config, always consult the official documentation for your software.

\section{Terminal}
Since we will be doing all of our changes with a text editor that lives inside the terminal, it makes sense to customize it first. 
We will be doing this by changing:
\begin{itemize}
    \item Cursor
    \item Colors
    \item Font
    \item Blur \& Transparency
    \item Scale
    \item Misc
\end{itemize}

Everything we will need is inside the configuration file for the terminal emulator, which should already be inside \textbf{~/.config/kitty/kitty.conf}

%CURSOR CONFIG
\subsection{Cursor}
Whoever uses the block style cursor is either very old and nostalgic, or a complete psychopath, so we need the I-Beam style.
%TODO insert image of block and beam here
To to this we need to change
\begin{lstlisting}
    cursor_shape block
\end{lstlisting}
To
\begin{lstlisting}
    cursor_shape beam
\end{lstlisting}

%COLOR CONFIG
\subsection{Color}
Now that we can actually use the terminal and edit text without dying, we can change the colors. We will be doing the "main" colors first, and the other colors after. So first we need the text color, the background color, and the cursor color.
This is really simple, it's just a matter of adding a hex code to a line:
\subsubsection{Main colors}
\begin{lstlisting}
        #foreground is the text color
   foreground #ffffff 
        #background is the terminal's background color
   background #000000
        #cursor is the cursor color, again, what a surpise!
   cursor #ffffff
\end{lstlisting}
Now that our text\footnote{I recommend setting the cursor color the same as the foreground, for visibility reasons, but you do not have to do so} is white and the background is black, we can move on to the other colors

\subsubsection{Secondary colors}
The terminal now is readable, but is very basic. Let's add some paint to it
Our terminal of choice has a 16 color palette, with 8 dull colors and 8 bright ones. You can customize the whole palette of the 256 terminal colors, but the basic ones are:
\begin{itemize}
        %TODO   check color codes
    \item Black
        \begin{lstlisting}
            color0 #000000
            color8 #767676
        \end{lstlisting}
    \item Red
        \begin{lstlisting}
            color1 #cc0403
            color9 #f2201f
        \end{lstlisting}
    \item Green
        \begin{lstlisting}
            color2 #cecb00
            color10 #fffd00
        \end{lstlisting}
    \item Yellow
        \begin{lstlisting}
            color3 #0d73cc
            color11 #1a8fff
        \end{lstlisting}
\pagebreak
    \item Blue
        \begin{lstlisting}
            color4 #cb1ed1
            color12 #1a8fff
        \end{lstlisting}
    \item Magenta
        \begin{lstlisting}
            color5 #0dcdc
            color13 #f228ff
        \end{lstlisting}
    \item Cyan
        \begin{lstlisting}
            color6 #0dcdcd
            color14 #14ffff
        \end{lstlisting}
    \item White
        \begin{lstlisting}
            color7 #dddddd
            color15 #ffffff
        \end{lstlisting}
\end{itemize}
These are the default ones, feel free to change them however you see fit

\subsection{Font}
Changing font is also really simple

\begin{lstlisting}
        #this is the actual font, in this case, the Ubuntu one
    font_familiy Adwaita Mono
        #these are the specific mods
    bold_font auto
    italic_font auto
    bold_italic_font auto
\end{lstlisting}

\subsection{Blur \& Transparency}
It's not a real Linux rice if your terminal is opaque. Unless you're going for that vibe, in which case, just skip ahead.
To change Transparency and Blur just change the homonimous rule:
\begin{lstlisting}
        # 1.0 means fully opaque, 0 means fully transparent
    background_opacity 0.8
        # 0 means no blur, increase to preference
    background_blur 4
\end{lstlisting}

\subsection{Scale}
You may need to have a bigger or smaller text inside the terminal. For this we will set a default one, and keybinds to increase and decrease it by intervals

\subsection{Misc}
